% Appendix B — Extended Discussion
% Contains discussion paragraphs moved from the main text during compression.

\section{Extended Discussion}\label{app:discussion}

%%% CF Dual Mechanism %%%

\paragraph{CF Dual Mechanism.}
CF violations arise from two distinct mechanisms:
(1)~\emph{Registry override}: the semantic constraint forces a domain-compatible type that conflicts with the entity's previously registered type.
The entity appears in multiple relations with different type requirements, and the per-relation mask selects the locally valid type without consulting prior assignments.
(2)~\emph{Grammar-first fallback} (\S\ref{sec:interaction}): when the semantic mask conflicts with all grammar-valid tokens, the semantic layer is dropped entirely, allowing the grammar layer to generate any structurally valid type, which may differ from the entity's registered type.

These mechanisms have sharply different model profiles.
Qwen3-8B triggers zero grammar-first fallbacks across all experiments, attributing 100\% of its CF violations to registry override.
Llama-3.1-8B, by contrast, triggers 5{,}882 fallbacks on SciERC (10.7 per document) and 57{,}825 on DocRED (63.8 per document), making grammar-first fallback the dominant CF source.
This asymmetry is consistent with the cross-model CF gap: Qwen achieves CF\,=\,2.57\% on DocRED vs.\ Llama's 12.02\%, despite both using identical constraint configurations.
The difference traces to Llama's less calibrated type predictions, which more frequently conflict with the semantic mask's valid set, triggering fallback.

%%% Fallback Behavior and Effective Coverage %%%

\paragraph{Fallback Behavior and Effective Coverage.}
\label{app:fallback-coverage}
The two-layer architecture produces two distinct fallback mechanisms with different semantic consequences.
\emph{Grammar-first fallback} occurs when $S(t) \cap G(t) = \emptyset$: the semantic layer deactivates as a safety valve, preserving structural validity at the cost of potential type violations.
The token-level fallback rate (36.6\% on SciERC Llama) is inflated by continuation tokens of multi-token type names (e.g., \textsc{OtherScientificTerm}, averaging 2.87 BPE tokens); at semantically meaningful first-token positions, the rate is 0.04--5.7\% (zero-shot) and drops to $\leq$\,0.11\% under fine-tuning.
These events account for the residual DR that persists after constraint enforcement.
\emph{Registry-consistent fallback} (4.5\%) occurs when the entity's registered type conflicts with domain/range constraints; the system defaults to the registered type for cross-triple consistency.
Together, these mechanisms produce the characteristic profile of composable constraint systems: residual DR traces to grammar-first fallback events, while CF violations arise from registry override events where per-relation enforcement assigns conflicting types.

We define \emph{effective semantic constraint coverage} as $1 - f_{\text{gf}}$, where $f_{\text{gf}}$ is the grammar-first fallback rate.
On SciERC, Llama-3.1-8B shows $f_{\text{gf}} = 36.6\%$ (6{,}237 of 17{,}024 mask calls), yielding 63.4\% effective coverage.
On DocRED, by contrast, Qwen3-8B shows $f_{\text{gf}} = 0.08\%$ (46 of 54{,}859), yielding 99.9\% coverage.
The large discrepancy is primarily a tokenization artifact: SciERC entity types are multi-token (e.g., \textsc{OtherScientificTerm}, averaging 2.87 mask calls per type field), and continuation tokens necessarily conflict with the semantic mask, inflating the per-mask-call fallback rate.
DocRED types (\textsc{PER}, \textsc{ORG}, \textsc{LOC}, \textsc{TIME}, \textsc{NUM}, \textsc{MISC}) are all single-token, eliminating this inflation.
This metric directly quantifies the layer interaction (\S\ref{sec:interaction}): it measures the fraction of decoding steps where both constraint layers are jointly active, providing a fine-grained view beyond aggregate violation rates.
Qwen3-8B, with zero grammar-first fallbacks, achieves 100\% effective coverage on SciERC, consistent with its lower residual DR (0.93\% vs.\ Llama's 0.08\%).

%%% Fallback Causal Tracing %%%

\paragraph{Fallback Causal Tracing.}
\label{app:causal-tracing}

To establish that grammar-first fallback is the \emph{sole} mechanism producing residual DR violations, we instrument the constraint processor to log, for each violation, whether the semantic mask was active ($S(t) \cap G(t) \neq \emptyset$) or bypassed via fallback ($S(t) \cap G(t) = \emptyset$) at the type-bearing position.
Table~\ref{tab:causal-tracing} reports the results across four SciERC configurations.

\begin{table}[t]
\centering
\small
\setlength{\tabcolsep}{4pt}
\resizebox{\columnwidth}{!}{%
\begin{tabular}{@{}ll rrr@{}}
\toprule
\textbf{Model} & \textbf{Setting} & \textbf{Fallback Rate} & \textbf{DR} & \textbf{From FB} \\
\midrule
Qwen3-8B  & FT S1 & 0/2256 (0\%)    & 0  & N/A \\
Qwen3-8B  & ZS S1 & 16/3458 (0.46\%) & 16 & 100\% \\
Llama-3.1-8B & FT S1 & 1/2162 (0.05\%) & 1  & 100\% \\
Llama-3.1-8B & ZS S1 & 1/2630 (0.04\%) & 1  & 100\% \\
\bottomrule
\end{tabular}}
\caption{Fallback causal tracing on SciERC (551 documents). Fallback Rate: type-bearing positions where $S(t) \cap G(t) = \emptyset$ triggers grammar-first fallback. DR: residual domain/range violations. From FB: fraction of DR violations occurring at fallback positions. All violations occur exclusively at tail\_type positions (range constraints are tighter than domain constraints).}
\label{tab:causal-tracing}
\end{table}

When the semantic mask is active, $\text{DR}_{\text{independent}} = 0$: the type constraint processor structurally prevents domain/range violations.
All 18 residual DR violations across four configurations occur exclusively at fallback positions, confirming that grammar-first fallback is a necessary condition for residual DR.
Fine-tuning reduces fallback to near-zero ($\leq$0.05\%), effectively eliminating the sole source of residual violations.
All fallback events occur at tail\_type positions, where $|\operatorname{range}(r, \tau^h)|$ is smaller than $|\operatorname{dom}(r)|$, making the intersection $S(t) \cap G(t)$ more likely to be empty.

%%% Enabling Dependency: Confound Analysis %%%

\paragraph{Enabling Dependency: Confound Analysis.}
\label{app:enabling-confound}

The scale-dependent layer interaction (\S\ref{sec:method}) shows that the semantic layer achieves only 24\% DR reduction without the grammar layer (Table~\ref{tab:scierc-ablation}).
A potential confound is that the semantic layer depends not only on grammar-enforced vocabulary membership but also on structural enforcement for positional awareness: the semantic layer identifies type-bearing positions by pattern-matching the decoded prefix against the JSON schema, which is reliable only when the grammar layer guarantees valid JSON structure.
Without grammar enforcement, malformed structure can cause the semantic layer to misidentify type-bearing positions, reducing effective coverage.
We quantify this via the semantic-only ablation (Table~\ref{tab:scierc-ablation}, row~3): CF rises from 0.08\% to 2.71\%, indicating that 2.63\,pp of additional violations trace to structural ambiguity.
However, the dominant enabling mechanism remains vocabulary membership: the semantic layer's domain/range lookup requires $r_j \in \mathcal{R}$, which only the grammar layer ensures; without it, OOV relation names produce null constraint lookups.
The confound reinforces rather than undermines this interaction: both vocabulary and structural contributions of the grammar layer are prerequisites for effective semantic enforcement.

%%% CF-DR Trade-off Analysis %%%

\paragraph{CF-DR Trade-off Analysis.}
\label{app:cf-dr-tradeoff}

CF violations are the most prominent residual limitation of OntoCoDec under zero-shot decoding (9.90\% on SciERC, up to 12.02\% on Llama DocRED).
We treat this honestly: per-relation independent enforcement cannot eliminate CF without trading it against DR.
Table~\ref{tab:cf-dr-tradeoff} reports a pilot ablation in which we replace the grammar-first fallback with a CF-priority fallback that forces the entity's registered type when it conflicts with $\sigma(r)$.

\begin{table}[t]
\centering
\small
\begin{tabular}{@{}l rrrr@{}}
\toprule
\textbf{Configuration} & \textbf{DR\,(\%)} & \textbf{CF\,(\%)} & \textbf{DR+CF} & \textbf{$f_{\text{gf}}$\,(\%)} \\
\midrule
S1 (baseline) & 0.93 & 9.90 & 10.83 & --- \\
CF fallback        & 14.78 & 0.04 & 14.82 & 36.64 \\
\midrule
$\Delta$ & $+$13.85 & $-$9.86 & $+$3.99 & --- \\
\bottomrule
\end{tabular}
\caption{CF-DR trade-off pilot on SciERC (Qwen3-8B, 551 documents). CF fallback eliminates CF but re-routes violations to DR, increasing total violations (DR+CF) by 3.99\,pp. $f_{\text{gf}}$: token-level grammar-first fallback rate at type-bearing positions (includes continuation tokens of multi-token type names; first-token position fallback rate is 0.46\% for SciERC Qwen; \S\ref{sec:greedy-limitation}).}
\label{tab:cf-dr-tradeoff}
\end{table}

The CF fallback variant nearly eliminates CF (9.90\%\,$\to$\,0.04\%, $-$9.86\,pp) but at the cost of a 15.9$\times$ increase in DR (0.93\%\,$\to$\,14.78\%, $+$13.85\,pp).
This is not a net improvement: total violations (DR+CF) rise from 10.83\% to 14.82\% ($+$3.99\,pp), so the change \emph{re-routes} violations rather than reducing them.
The mechanism is mechanical: each token position with conflicting masks must commit to one violation type; when the CF fallback forces the system to emit the entity's registered type even though it lies outside $\sigma(r)$, the resulting triple satisfies cross-entity consistency but violates domain/range.
The token-level grammar-first fallback is triggered at 36.64\% of type-bearing positions (0.46\% at first-token positions; \S\ref{sec:greedy-limitation}); this conflict is structural, not a tunable hyperparameter.

We interpret this as a fundamental tension in composable constraint systems with \emph{per-relation independent} enforcement: enforcement priority can be allocated between DR and CF but cannot independently minimize both.
Within this design, the current S1 configuration sits near the Pareto frontier of the trade-off, achieving 0.93\% DR at the cost of 9.90\% CF.
The direction of the trade-off is consistent beyond SciERC/Qwen3-8B: across all six zero-shot configurations (3 models $\times$ 2 datasets), DR and CF are descriptively inversely associated, with no setting showing both high DR and high CF simultaneously (Table~\ref{tab:cross-model}); however, $n{=}6$ is insufficient for statistical inference on this correlation.
The pattern is a structural consequence of per-relation enforcement with grammar-first fallback, not a model- or dataset-specific artifact; we select SciERC/Qwen3-8B for detailed ablation because it exhibits the highest CF (9.90\%), making the trade-off most visible.
While the inverse association holds across all configurations, the absolute trade-off magnitude may differ in settings with higher baseline CF, such as Llama DocRED (CF = 12.02\%), where the larger CF headroom could yield proportionally larger shifts under alternative fallback priorities.
The CF residual is concentrated in schema ambiguity rather than enforcement failure: 63.3\% of CF violations on SciERC involve the catch-all \textsc{Generic} type (Table~\ref{tab:cf-distribution}), which competes with every content type depending on anaphoric vs.\ referential use, a distinction the per-slot decoder cannot resolve without document-level coreference.

Two complementary paths can lower CF without re-routing it to DR.
(i)~\emph{Reduce the baseline conflict rate.} Composing OntoCoDec with LoRA fine-tuning on SciERC reduces CF by 88\% relative to the zero-shot setting (9.90\%\,$\to$\,1.18\%; Table~\ref{tab:composability}), because the fine-tuned model's per-slot type predictions agree more often with the entity registry, so the per-relation mask and the registry intersect more frequently.
The remaining 1.18\% still traces to \textsc{Generic}$\leftrightarrow$X schema ambiguity, so this path is empirically validated on one dataset and one fine-tuning method but does not generalize without further evidence.
(ii)~\emph{Replace per-relation enforcement with global entity-level type consensus during decoding.} This is outside the scope of OntoCoDec's forward-only autoregressive design and is left to future work.

%%% Dual-Factor Quantitative Analysis %%%

\paragraph{Dual-Factor Quantitative Analysis.}
\label{app:dual-factor}

Figure~\ref{fig:dual-factor-scatter} plots residual DR against the product $\bar{\kappa} \times \text{RV}_{\text{baseline}}$ for model$\times$dataset configurations.
For \emph{calibrated} models (Qwen3-8B, Llama-3.1-8B; unconstrained ET${<}$3\% on both datasets), a linear fit yields $y = 0.366x - 0.138$ with $R^2 = 0.886$ ($n{=}4$, $F(1,2) = 15.54$, $p \approx 0.058$); with only four data points (df $= 2$) and $p > 0.05$, this constitutes a descriptive characterization rather than a confirmatory statistical test.
The strong fit supports the dual-factor pattern from \S\ref{sec:composition}: residual DR is governed by the product of schema tightness and baseline violation rate.
Mistral-7B (unconstrained ET${=}$7.58\% on DocRED) deviates from this pattern: its residual DR on DocRED (33.21\%) substantially exceeds the linear prediction. When model calibration is insufficient (ET${>}$5\% on either dataset), residual DR is dominated by entity misclassification rather than constraint-search failure.
This motivates a principled exclusion criterion: models with unconstrained ET${>}$5\% on either dataset fall outside the regime where the tightness$\times$RV product is the primary DR driver.
The two DocRED points (high $\bar{\kappa}{=}0.842$) separate cleanly from the two SciERC points (low $\bar{\kappa}{=}0.333$), and within each dataset the model with higher baseline RV shows higher residual DR.

% Dual-factor scatter moved to main text (Figure~\ref{fig:dual-factor-scatter}).

%%% Dataset-Dependent F1 Impact (moved from §7) %%%

\paragraph{Dataset-Dependent F1 Impact.}
The divergence in F1 significance between DocRED ($+$0.49\,pp, $p < 0.001$) and SciERC ($+$0.13\,pp, $p = 0.58$) traces to a single mechanism: grammar-layer elimination of out-of-vocabulary relation names.
DocRED's baseline RV of 29.66\% produces 5{,}189 OOV predictions counted as false positives; eliminating them yields a significant precision gain.
SciERC's baseline RV is only 1.39\% (24 OOV predictions), too few to produce a detectable F1 effect.
Llama-3.1-8B's lower baseline RV (16.03\%) yields a smaller but directionally consistent improvement ($+$0.29\,pp vs.\ $+$0.49\,pp); F1 gain magnitude is proportional to baseline OOV rate.

%%% Bootstrap CIs and Effect Sizes %%%

\paragraph{Bootstrap Confidence Intervals and Effect Sizes.}
\label{app:bootstrap-effect-sizes}

Table~\ref{tab:bootstrap-effect-sizes} reports per-configuration bootstrap CIs and Cohen's $d$ effect sizes (10{,}000 iterations, document-level resampling, seed=42) for the fine-tuning composability results (Table~\ref{tab:composability}).
After Bonferroni correction for $k{=}4$ comparisons ($\alpha_{\text{adj}}{=}0.0125$), no R-F1 difference reaches significance; all R-F1 effect sizes are negligible to small ($|d| \leq 0.37$; Cohen, 1988).
Per-configuration corrected $p$-values: Qwen3-8B/SciERC $p_B{=}1.00$, Qwen3-8B/DocRED $p_B{=}0.224$, Llama-3.1-8B/SciERC $p_B{=}1.00$, Llama-3.1-8B/DocRED $p_B{=}1.00$.
Constraint enforcement thus preserves extraction quality while achieving large, highly significant reductions in domain-range violations ($|d| \geq 3.0$, all $p_{\text{Bonf}} < 0.005$).

\begin{table}[t]
\centering
\small
\setlength{\tabcolsep}{3pt}
\resizebox{\columnwidth}{!}{%
\begin{tabular}{@{}ll rrrr@{}}
\toprule
\textbf{Model} & \textbf{Dataset} & $\Delta$\textbf{(pp)} & \textbf{95\% CI} & $p$ & $d$ \\
\midrule
\multicolumn{6}{@{}l}{\textbf{R-F1 difference (S1 $-$ unconstrained):}} \\
Qwen3-8B & SciERC & $-$0.10 & [$-$0.30, 0.00] & .726 & $-$0.05 \\
Qwen3-8B & DocRED & $-$0.31 & [$-$0.69, $+$0.01] & .056 & $-$0.37 \\
Llama-3.1-8B & SciERC & $+$0.02 & [0.00, $+$0.07] & .736 & $+$0.01 \\
Llama-3.1-8B & DocRED & $+$0.06 & [$-$0.11, $+$0.25] & .535 & $+$0.07 \\
\midrule
\multicolumn{6}{@{}l}{\textbf{DR\% reduction:}} \\
Qwen3-8B & SciERC & $-$1.95 & [$-$2.87, $-$1.11] & $<$.001 & $-$6.15 \\
Qwen3-8B & DocRED & $-$0.75 & [$-$1.07, $-$0.48] & $<$.001 & $-$7.00 \\
Llama-3.1-8B & SciERC & $-$1.02 & [$-$1.98, $-$0.29] & $<$.001 & $-$3.17 \\
Llama-3.1-8B & DocRED & $-$0.29 & [$-$0.53, $-$0.10] & $<$.001 & $-$3.02 \\
\bottomrule
\end{tabular}}
\caption{Bootstrap CIs and Cohen's $d$ for Table~\ref{tab:composability} (10{,}000 iterations, document-level resampling, seed=42).
$d$: Cohen's $d$ from pooled bootstrap SD.
Top: R-F1 (S1 $-$ unconstrained); no comparison significant after Bonferroni correction ($k{=}4$, $\alpha_{\text{adj}}{=}0.0125$); $|d| \leq 0.37$ (negligible--small).
Bottom: DR\% reduction; all $p_{\text{Bonf}} < 0.005$, $|d| \geq 3.0$ (large).}
\label{tab:bootstrap-effect-sizes}
\end{table}

%%% Paragraph 1: Low F1 Practical Significance %%%

\paragraph{Practical Significance of Zero-Shot F1.}
Our reported relation F1 of approximately 8\% on both SciERC and DocRED reflects the inherent difficulty of zero-shot closed IE, where the model receives no task-specific training examples.
For reference, supervised state-of-the-art systems on SciERC achieve approximately 49--53\% relation F1 with full training data~\citep{zhong2021frustratingly,ye2022packed}.
Constraint enforcement operates largely independently of the extraction model's intrinsic capability: it restricts the output space to ontologically valid triples regardless of the underlying model's accuracy.
Even after LoRA fine-tuning, residual DR of 1--2\% persists on tight-constraint relations (Table~\ref{tab:composability}); constraints remain necessary at deployment-grade quality.
In our evaluation, F1 serves as a \emph{control variable} verifying that constraint enforcement does not degrade extraction quality, not as an optimization target.
The contribution of this work is the constraint enforcement method itself; the zero-shot setting stress-tests constraint robustness under maximum model uncertainty, while composability with fine-tuning (\S\ref{sec:composability}) validates practical deployment value.

%%% Paragraph 4: Decoding-Time vs. Post-Hoc Enforcement %%%

\paragraph{Decoding-Time vs.\ Post-Hoc Enforcement.}
Post-hoc filtering (generating triples without constraints, then removing schema-violating predictions) yields higher precision and consequently higher F1 than OntoCoDec on both datasets (SciERC 8.63 vs.\ 8.36; DocRED 9.61 vs.\ 8.23).
% posthoc: SciERC F1=8.63, DocRED F1=9.61; OntoCoDec: SciERC R-F1=8.36, DocRED R-F1=8.23
However, this precision gain comes at substantial cost.
On DocRED, post-hoc filtering discards 48.1\% of all predicted triples, including 127 true positives (11.2\% of correct predictions); recall drops from 9.46\% to 8.40\%.
% posthoc_docred: rejection=48.1%, tp_lost=127/1137, R: 9.46→8.40
Entity type violations remain unresolved (ET\,=\,2.47\% vs.\ OntoCoDec's 0\%), as filtering targets only domain/range and vocabulary violations.
OntoCoDec's decoding-time approach avoids this information loss by guiding generation toward compliant outputs rather than discarding non-compliant ones, preserving both recall and entity type validity.
The two strategies are complementary: post-hoc filtering can serve as an additional safety net for residual violations that decoding-time enforcement does not fully resolve.

%%% CF Violation Distribution %%%

\paragraph{Cross-Field Violation Distribution.}
\label{app:cf-distribution}

Table~\ref{tab:cf-distribution} reports the complete distribution of conflicting-type pairs for SciERC zero-shot S1 (248 CF violations across 245 unique entity-document pairs; no mention triggered $\geq$3 types).
Conflicts involving the \textsc{Generic} type account for 63.3\% (155/245) of all violations.
\textsc{Generic} is SciERC's anaphoric/placeholder label (``this approach'', ``the method''), which competes with every content type depending on whether the mention is used anaphorically or referentially, a distinction the per-slot constrained decoder cannot resolve without document-level coreference.

The pattern generalizes across models and datasets: on Llama-3.1-8B SciERC, \textsc{Generic}$\leftrightarrow$\textsc{OtherScientificTerm} remains the top pair (53/162, 32.7\%); on DocRED, the analogous catch-all type \textsc{MISC} drives 76.6\% of Qwen's CF violations.
Under unconstrained decoding, CF rates collapse to $\leq$1\% on both datasets; CF violations are a side-effect of per-slot type commitment rather than a model deficiency.

\begin{table}[t]
\centering
\small
\begin{tabular}{@{}lr r@{}}
\toprule
\textbf{Conflicting Type Pair} & \textbf{Count} & \textbf{\%} \\
\midrule
Generic $\leftrightarrow$ OtherScientificTerm & 85 & 34.7 \\
Generic $\leftrightarrow$ Material & 35 & 14.3 \\
Generic $\leftrightarrow$ Metric & 30 & 12.2 \\
Method $\leftrightarrow$ Task & 20 & 8.2 \\
Material $\leftrightarrow$ Method & 17 & 6.9 \\
OtherScientificTerm $\leftrightarrow$ Task & 12 & 4.9 \\
\midrule
10 further pairs (each $\leq$3.3\%) & 46 & 18.8 \\
\midrule
\textbf{Total} & \textbf{245} & 100.0 \\
\bottomrule
\end{tabular}
\caption{CF violation distribution on SciERC zero-shot S1 (Qwen3-8B, 551 documents).
All 245 unique entity-document pairs involve pairwise conflicts.
Conflicts involving \textsc{Generic} account for 63.3\%.}
\label{tab:cf-distribution}
\end{table}

%%% Cascade Error Analysis %%%

\paragraph{Cascade Error Analysis.}
\label{app:cascade-analysis}

Table~\ref{tab:cascade-breakdown} decomposes zero-shot CF violations by source mechanism.
The dominant source is hallucinated entities (58.1\%): the model generates entity names absent from the input document, and the entity registry assigns a type based on the first occurrence that may conflict with a subsequent valid entity sharing the same surface form.
True cascade errors---where first-seen type registration for a \emph{real} entity locks an incorrect type that blocks a subsequent valid assignment---account for only 7.3\% (18/248).
Correct type locks (first-seen type is correct and prevents an incorrect subsequent assignment) are 2.5$\times$ more frequent (18.1\%, 45/248), indicating that the first-seen registration mechanism provides a net benefit for type consistency.

Under fine-tuning (22 CF violations, 20 documents), cascade errors are completely eliminated (0/22): LoRA training calibrates entity typing sufficiently that first-seen registrations are consistently correct.
This confirms that cascade errors are a model calibration issue exposed by the constraint mechanism, not an inherent deficiency of first-seen registration.

\begin{table}[h]
\centering
\small
\begin{tabular}{@{}l rr@{}}
\toprule
\textbf{CF Source Category} & \textbf{Count} & \textbf{\%} \\
\midrule
Hallucinated entity & 144 & 58.1 \\
Correct type lock & 45 & 18.1 \\
True cascade error & 18 & 7.3 \\
Other / ambiguous & 41 & 16.5 \\
\midrule
\textbf{Total} & \textbf{248} & 100.0 \\
\bottomrule
\end{tabular}
\caption{CF violation source breakdown on SciERC zero-shot S1 (Qwen3-8B, 551 documents).
``Hallucinated entity'': generated entity absent from input.
``True cascade error'': first-seen lock on incorrect type blocks valid subsequent assignment.
``Correct type lock'': first-seen type correct, prevents incorrect subsequent type.
Under fine-tuning, cascade errors are completely eliminated (0/22 CF violations).}
\label{tab:cascade-breakdown}
\end{table}

%%% Per-Relation DR Elimination under Fine-Tuning %%%

\paragraph{Per-Relation DR Elimination under Fine-Tuning.}
\label{app:lora-per-relation}

Table~\ref{tab:lora-per-relation} presents the per-relation breakdown of DR elimination when OntoCoDec is composed with LoRA fine-tuning on SciERC.
Fine-tuning alone eliminates DR on four of seven relations but leaves non-zero violations on the three tightest-constraint relations: \textsc{Evaluate-for} ($\kappa{=}0.361$, 13.13\%), \textsc{Part-of} ($\kappa{=}0.417$, 5.77\%), and \textsc{Compare} ($\kappa{=}0.528$, 5.00\%).
Adding OntoCoDec constraints eliminates all remaining violations across all relations.
One \textsc{Evaluate-for} triple (99\,$\to$\,98) is redirected to \textsc{Used-for} (634\,$\to$\,635) by token-level masking: the head entity's registered type falls outside \textsc{Evaluate-for}'s domain, and the semantic mask steers generation to the nearest valid relation.

\begin{table}[t]
\centering
\small
\begin{tabular}{@{}l r rr rr@{}}
\toprule
& & \multicolumn{2}{c}{\textbf{LoRA Unconstr.}} & \multicolumn{2}{c}{\textbf{LoRA + S1}} \\
\cmidrule(lr){3-4} \cmidrule(l){5-6}
\textbf{Relation} & $\kappa$ & DR\,(\%) & $n$ & DR\,(\%) & $n$ \\
\midrule
Evaluate-for & 0.361 & 13.13 & 99 & \textbf{0} & 98 \\
Part-of      & 0.417 & 5.77  & 52 & \textbf{0} & 52 \\
Compare      & 0.528 & 5.00  & 40 & \textbf{0} & 40 \\
Conjunction  & 0.306 & 1.31  & 153 & \textbf{0} & 153 \\
Feature-of   & 0.278 & 1.35  & 74 & \textbf{0} & 74 \\
Used-for     & 0.028 & 0.16  & 634 & \textbf{0} & 635 \\
Hyponym-of   & 0.417 & 0     & --- & \textbf{0} & --- \\
\bottomrule
\end{tabular}
\caption{Per-relation DR elimination on SciERC with LoRA fine-tuning (Qwen3-8B).
$\kappa$: constraint tightness; $n$: predicted triple count.
Used-for gains one triple from Evaluate-for redirection.}
\label{tab:lora-per-relation}
\end{table}

%%% Fine-Tuned Four-Way Ablation %%%

\paragraph{Fine-Tuned Four-Way Ablation.}
\label{app:ft-ablation}

Table~\ref{tab:ft-ablation} presents the four-way constraint ablation on fine-tuned models across both datasets, isolating each layer's contribution when the base model has already internalized schema knowledge.

\begin{table}[t]
\centering
\small
\setlength{\tabcolsep}{3pt}
\resizebox{\columnwidth}{!}{%
\begin{tabular}{@{}lll rrrrr rr@{}}
\toprule
\textbf{Dataset} & \textbf{Model} & \textbf{Mode} & \textbf{DR\,$\downarrow$} & \textbf{RV\,$\downarrow$} & \textbf{ET\,$\downarrow$} & \textbf{CF} & \textbf{PF} & \textbf{E-F1\,$\uparrow$} & \textbf{R-F1\,$\uparrow$} \\
\midrule
\multirow{8}{*}{SciERC}
  & \multirow{4}{*}{Qwen3-8B}
    & Unconstrained & 1.95 & 0.00 & 0.00 & 0.00 & 0.00 & 67.05 & 43.86 \\
  & & Grammar       & 1.90 & 0.00 & 0.00 & 0.00 & 0.00 & 67.05 & 43.86 \\
  & & Semantic      & 0.60 & 0.00 & 0.00 & 0.18 & 0.00 & 67.05 & 43.69 \\
  & & S1       & \textbf{0}    & 0.00 & 0.00 & 1.18 & 0.00 & 67.05 & 43.77 \\
  \cmidrule(l){2-10}
  & \multirow{4}{*}{Llama-3.1-8B}
    & Unconstrained & 1.11 & 0.00 & 0.00 & 0.00 & 0.00 & 67.76 & 45.45 \\
  & & Grammar       & 1.11 & 0.00 & 0.00 & 0.00 & 0.00 & 67.76 & 45.45 \\
  & & Semantic      & 0.65 & 0.00 & 0.00 & 0.10 & 0.00 & 67.76 & 45.47 \\
  & & S1       & \textbf{0.09} & 0.00 & 0.00 & 0.46 & 0.00 & 67.76 & 45.47 \\
\midrule
\multirow{8}{*}{DocRED}
  & \multirow{4}{*}{Qwen3-8B}
    & Unconstrained & 0.76 & 0.27 & 0.00 & 0.12 & 7.11 & 82.15 & 40.86 \\
  & & Grammar       & 0.61 & 0.00 & 0.00 & 0.11 & 7.01 & 82.19 & 40.85 \\
  & & Semantic      & 0.19 & 0.23 & 0.00 & 0.27 & 7.52 & 82.15 & 40.85 \\
  & & S1       & \textbf{0.01} & 0.00 & 0.00 & 0.41 & 8.52 & 81.43 & 40.56 \\
  \cmidrule(l){2-10}
  & \multirow{4}{*}{Llama-3.1-8B}
    & Unconstrained & 0.49 & 0.15 & 0.00 & 0.08 & 2.30 & 83.93 & 43.72 \\
  & & Grammar       & 0.51 & 0.00 & 0.00 & 0.12 & 2.30 & 84.04 & 43.72 \\
  & & Semantic      & 0.48 & 0.15 & 0.00 & 0.16 & 2.40 & 84.02 & 43.85 \\
  & & S1       & \textbf{0.21} & 0.00 & 0.00 & 0.21 & 2.51 & 83.94 & 43.78 \\
\bottomrule
\end{tabular}}%
\caption{Four-way constraint ablation on fine-tuned models (LoRA, 2 models $\times$ 2 datasets).
DR/RV/ET: violation rates (\%) over predicted relations; CF: over predicted entities; PF: parse failure rate (\%); F1 $\times$100.
Grammar layer is a no-op post-fine-tuning (DR change $<$0.2\,pp in all 4 settings; RV and ET already 0\% without constraints).
Semantic layer is the sole effective DR defense (41--69\% on SciERC, 2--75\% on DocRED); the near-zero Llama DocRED reduction reflects a near-floor baseline (DR${=}$0.49\%).
F1 cost is negligible: max $\Delta$R-F1 = $-$0.30\,pp (Qwen DocRED, full vs.\ unconstrained).}
\label{tab:ft-ablation}
\end{table}

The grammar layer is effectively a no-op for fine-tuned models across all four settings: RV and ET are already 0\% without constraints, and grammar-only DR change is $<$0.2\,pp (4/4 settings).
The semantic layer is the sole effective DR defense under fine-tuning, independently reducing DR by 41--69\% on SciERC (Qwen $-$69\%, Llama $-$41\%) and 2--75\% on DocRED (Qwen $-$75\%, Llama $-$2\%); the near-zero Llama DocRED reduction reflects a near-floor baseline (DR${=}$0.49\%, 80 violations in 16{,}074 triples), where the semantic layer has minimal room to improve.
Adding both layers achieves near-complete elimination (DR\,$\leq$\,0.21\% in all 4 settings): even after fine-tuning, the semantic layer's pair-based enforcement captures violations that training alone cannot eliminate, specifically on tight-constraint relations where the valid type-pair space is small.
F1 cost is negligible: the maximum R-F1 decrease across S1 vs.\ unconstrained is $-$0.30\,pp (Qwen DocRED).

%%% DocRED Zero-Shot Ablation (moved from main text) %%%

\paragraph{DocRED Zero-Shot Ablation.}
\label{app:docred-zs-ablation}

Table~\ref{tab:docred-ablation} presents the four-way constraint ablation on DocRED under zero-shot decoding, parallel to the SciERC ablation in Table~\ref{tab:scierc-ablation}.

\begin{table}[t]
\centering
\small
\setlength{\tabcolsep}{3.6pt}
\resizebox{\columnwidth}{!}{%
\begin{tabular}{@{}l rrrrrr rr@{}}
\toprule
& \multicolumn{6}{c}{\textbf{Violation Rates (\%)}}
& \multicolumn{2}{c}{\textbf{F1}} \\
\cmidrule(lr){2-7} \cmidrule(l){8-9}
\textbf{Mode}
  & DR\,$\downarrow$ & $\Delta$DR
  & RV\,$\downarrow$ & ET\,$\downarrow$ & CF & Tot.\,$\downarrow$
  & Rel.\,$\uparrow$ & Ent.\,$\uparrow$ \\
\midrule
Unconstrained
  & 18.46 & ---
  & 29.66 & 2.48
  & 0.10  & 18.56
  & 7.74  & 58.15 \\
Grammar
  & 26.17 & $\uparrow$42\%
  & \textbf{0.03} & \textbf{0}
  & 0.07  & 26.24
  & \textbf{8.31}  & \textbf{59.38} \\
Semantic
  & 13.84 & $\downarrow$25\%
  & 29.87 & 4.22
  & 1.05  & 14.89
  & 7.70  & 57.81 \\
S1
  & \textbf{10.21} & $\downarrow$\textbf{45\%}
  & \textbf{0.03} & \textbf{0}
  & 2.57  & \textbf{12.78}
  & 8.23  & 59.37 \\
\bottomrule
\end{tabular}}%
\caption{Four-way ablation on DocRED (998 documents, Qwen3-8B, zero-shot).
$\Delta$DR: relative change vs.\ unconstrained.
Tot.\ = DR + CF (total type-violation burden).
Best per column in \textbf{bold} (CF excluded; Tot.\ included).
F1 $\times$100. Bootstrap significance: R-F1 $+$0.49\,pp ($p < 0.001$).
Grammar-only $\uparrow$DR: vocabulary enforcement makes OOV outputs evaluable, exposing latent violations (see \S\ref{sec:docred-results}).}
\label{tab:docred-ablation}
\end{table}

%%% Residual DR Error Analysis (Fine-Tuned) %%%

\paragraph{Residual DR Error Analysis.}
\label{app:residual-dr-analysis}

Across all fine-tuned settings (2 models $\times$ 2 datasets $\times$ S1), only 36 domain/range violations remain.
Manual inspection reveals four root causes:

\begin{enumerate}[nosep,leftmargin=*]
\item \emph{Entity type errors} (39\%): the LLM assigns an incorrect entity type (e.g., labeling a geographic entity as \textsc{LOC} when the gold standard requires \textsc{MISC}), causing the semantic mask to permit a domain/range-incompatible relation.
\item \emph{Schema coverage gaps} (33\%): the frequency-based constraint table, constructed from training-set co-occurrence statistics, does not cover long-tail type pairs that are ontologically valid but rare. These violations are irreducible by decoding-time constraints alone.
\item \emph{Incorrect relation selection} (17\%): the model selects a semantically inappropriate relation whose domain/range constraints are incompatible with the correctly typed entities.
\item \emph{Direction reversal} (11\%): the model assigns the correct relation but inverts head/tail roles, placing entities in positions that violate the relation's asymmetric domain/range specification.
\end{enumerate}

Two dataset-specific patterns emerge.
Outside Llama DocRED, all residual violations (3 of 36, spanning Llama SciERC and Qwen DocRED) trace exclusively to schema coverage gaps: the model's entity typing is accurate, and the constraint system has reached its theoretical ceiling given the frequency-based constraint table.
On DocRED, Llama-3.1-8B contributes 33 of 36 residual violations (92\%), with entity type errors as the dominant cause (14/33, 42\%).
This asymmetry reflects Llama's weaker fine-tuned entity typing on DocRED's larger ontology (96 relations, 6 entity types) compared to Qwen, which achieves near-zero residual.
The implication is twofold: (i)~the irreducible floor from schema gaps is small ($\sim$3 violations across 4 settings), so the constraint system operates near its theoretical optimum; (ii)~residual violations from entity typing errors are addressable through stronger base models or improved entity typing modules, without modifying the constraint architecture.

%%% Post-Hoc Filtering Comparison (Extended) %%%

\paragraph{Post-Hoc Filtering: Extended Comparison.}
\label{app:posthoc-extended}

Table~\ref{tab:posthoc-extended} presents the full comparison across all six zero-shot and four fine-tuned configurations.

\begin{table}[t]
\centering
\small
\setlength{\tabcolsep}{3pt}
\resizebox{\columnwidth}{!}{%
\begin{tabular}{@{}lll rrr l@{}}
\toprule
\textbf{Setting} & \textbf{Model} & \textbf{Dataset} & \textbf{Reject\%} & \textbf{R-F1$_{\text{ph}}$} & \textbf{R-F1$_{\text{cd}}$} & \textbf{Winner} \\
\midrule
\multirow{6}{*}{ZS}
  & Qwen    & SciERC & 16.8 & 8.63  & 8.36  & ph \\
  & Llama   & SciERC & 13.9 & 2.71  & 2.54  & ph \\
  & Mistral & SciERC & 20.3 & 3.94  & 3.67  & ph \\
  & Qwen    & DocRED & 48.1 & 9.61  & 8.23  & ph \\
  & Llama   & DocRED & 42.2 & 7.44  & 6.27  & ph \\
  & Mistral & DocRED & \textbf{73.7} & 3.85  & \textbf{5.52} & cd \\
\midrule
\multirow{4}{*}{FT}
  & Qwen    & SciERC & 1.95 & 43.94 & 43.77 & ph \\
  & Llama   & SciERC & 1.11 & 45.23 & 45.47 & cd \\
  & Qwen    & DocRED$^\ddagger$ & 1.03 & 41.11 & 40.79 & ph \\
  & Llama   & DocRED & 0.64 & 43.92 & 43.78 & cd \\
\bottomrule
\end{tabular}}%
\caption{Post-hoc filtering (ph) vs.\ constrained decoding (cd).
In zero-shot settings, post-hoc F1 is higher in 5/6 groups (precision-dominated), but constrained decoding preserves higher recall in all 6/6 groups.
At reject rate $>$70\% (Mistral DocRED), post-hoc F1 collapses below constrained.
Under fine-tuning, rejection rates drop to 0.6--2.0\% and the two strategies are equivalent ($\Delta \leq 0.55$\,pp; each wins two groups).
F1 $\times$100.
$^\ddagger$Uses 4096-context LoRA variant; R-F1 differs from Table~\ref{tab:composability}'s standard LoRA configuration.}
\label{tab:posthoc-extended}
\end{table}

The crossover point occurs when post-hoc rejection exceeds $\sim$70\%: Mistral DocRED rejects 73.7\% of predicted triples, discarding enough true positives that constrained decoding's recall advantage dominates ($+$1.67\,pp).
OntoCoDec's advantage thus scales with violation severity: precisely when constraints are most needed, post-hoc filtering is most destructive.
Under fine-tuning, both strategies converge because few violations remain to filter or constrain ($\Delta \leq 0.55$\,pp across all four FT groups, with each strategy winning two).

%%% Zero-Shot Runtime Profiling %%%

\paragraph{Zero-Shot Runtime Profiling.}
\label{app:runtime}

On SciERC (Qwen3-8B, zero-shot greedy, NVIDIA RTX PRO 6000 Blackwell Server Edition), unconstrained inference averages 4.163\,s/doc and S1 constrained decoding averages 4.710\,s/doc ($+$13.1\% overhead).
On DocRED (same configuration), relative overhead is lower due to amortization: the fixed per-token constraint cost is spread over longer generation sequences.
XGrammar automaton compilation takes 0.044\,s; the full constraint pipeline (schema parsing, JSON Schema generation, grammar compilation) takes 0.33\,s on SciERC (Table~\ref{tab:scalability}), amortized to zero over the evaluation set.

\paragraph{Fine-Tuned Latency Profiling.}
\label{app:ft-latency}
Under LoRA fine-tuning (same hardware and dataset), unconstrained inference averages 2.612\,s/doc and S1 constrained decoding averages 2.873\,s/doc ($+$10.0\% overhead, $n{=}551$).
Fine-tuned models generate shorter, more structured outputs than zero-shot models (2.612\,s vs.\ 4.163\,s/doc unconstrained); absolute generation time drops by 37\%.
The relative overhead ($+$10.0\%) is slightly lower than zero-shot ($+$13.1\%), consistent with shorter sequences requiring fewer constraint checks per document.

%%% Fuzzy Relation Matching %%%

\paragraph{Relation Matching Protocol.}
\label{app:fuzzy-matching}
Relation F1 evaluation uses fuzzy normalization to handle model-specific formatting variants (e.g., underscores vs.\ spaces, abbreviations).
A predicted relation $r_p$ matches gold relation $r_g$ if: (1)~exact string match after case-insensitive normalization, or (2)~token overlap $>$80\% \emph{and} edit distance $\leq 2$.
This protocol affects primarily Mistral, which produces underscore-separated variants (e.g., ``head\_of\_government'' for ``head of government'').
On Mistral DocRED, fuzzy matching contributes 98 additional TPs ($+$0.84\,pp R-F1; strict R-F1\,=\,3.50\% vs.\ fuzzy 4.34\%).
For Qwen and Llama under grammar constraints, the grammar layer enforces exact relation names, making fuzzy matching irrelevant (0 additional matches).
All F1 figures in the main text use fuzzy normalization consistently across all models and conditions.

%%% Per-Relation Violation Landscape %%%

\paragraph{Per-Relation Violation Landscape.}
\label{app:per-relation-landscape}

Table~\ref{tab:prompt_comparison} compares violation rates under unconstrained and schema-informed prompting.
Figure~\ref{fig:per_relation_dr} shows the per-relation DR distribution on SciERC.

\begin{table}[t]
\centering
\small
\caption{Violation rates (\%) under unconstrained and schema-informed prompting (Qwen3-8B, zero-shot).}
\label{tab:prompt_comparison}
\resizebox{\columnwidth}{!}{%
\begin{tabular}{ll cccc}
\toprule
Dataset & Prompt & DR$\downarrow$ & RV$\downarrow$ & ET$\downarrow$ & CF$\downarrow$ \\
\midrule
\multirow{2}{*}{SciERC} & Unconstrained & 15.4 & 1.4 & 0.4 & 0.1 \\
 & Schema-informed & 5.8 & 0.3 & \textbf{0.0} & --- \\
\midrule
\multirow{2}{*}{DocRED} & Unconstrained & 18.5 & 29.7 & 2.5 & 0.1 \\
 & Schema-informed & 12.8 & 18.4 & 0.2 & 0.2 \\
\bottomrule
\end{tabular}}
\end{table}

%%% WikiEvents Extended %%%

\paragraph{WikiEvents Dataset Details.}
\label{app:wikievents}
WikiEvents~\cite{li2021document} is a document-level event argument extraction benchmark grounded in the DARPA KAIROS ontology.
It defines 56 event types, 80 argument roles, and 27 entity types, with 245 (event\_type, role)$\to$allowed\_entity\_types constraints.
The grammar layer targets EV (event vocabulary), RoV (role vocabulary), and ET; the semantic layer enforces AC (argument constraint) via role-conditioned masking, analogous to RE domain/range enforcement.

\paragraph{Constraint Tightness in WikiEvents.}
\label{app:wikievents-kappa}
WikiEvents role-level tightness (with $|\mathcal{E}|{=}27$) has mean $\bar{\kappa}{=}0.82$ and median 0.89; 84\% of roles have $\kappa \geq 0.80$ (4 or fewer valid entity types out of 27).
The tightness-violation pattern extends to event extraction: roles with $\kappa \geq 0.90$ exhibit the highest AC rate (50\%), compared to 21\% for medium-tightness ($0.80 \leq \kappa < 0.90$) and 30\% for loose ($\kappa < 0.80$) roles; this is consistent with the RE observation that constraint tightness correlates with violation frequency.
However, the correlation is modulated by entity type prototypicality: semantically intuitive roles (e.g., Transportation.Vehicle$\to$VEH) show 0\% AC despite $\kappa{=}0.963$, while roles that require rare KAIROS-specific types (e.g., Crime$\to$CRM, Sentence$\to$SEN) show near-100\% AC because LLMs default to common types (PER, ORG).
The loose group's elevated rate (30\%) is driven by high-frequency roles such as Attack.Target (126 args, 34\% AC, $\kappa{=}0.59$), where systematic model errors outweigh the permissive constraint space.

%%% Lucky Violation Analysis %%%

\paragraph{Lucky Violation Analysis.}
\label{app:lucky-violations}
Table~\ref{tab:lucky-violations} quantifies the recall cost of DR elimination.
A ``lucky violation'' is a constraint-removed triple that happens to match a gold label---an ontologically invalid prediction that is nonetheless correct.
Across four fine-tuned configurations, lucky violations average 0.24\% of gold triples (range 0.01--0.51\%), confirming that nearly all removed DR violations are false positives.
On Llama SciERC, all five lucky violations are recovered as legal type paths under constrained decoding (zero churn): the model finds an ontology-compliant route to the same prediction.

\begin{table}[h]
\centering
\small
\resizebox{\columnwidth}{!}{%
\begin{tabular}{@{}llrrrr@{}}
\toprule
\textbf{Model} & \textbf{Dataset} & \textbf{DR (\%)} & \textbf{Lucky} & \textbf{$\Delta$Recall} & \textbf{Churn} \\
& & \textbf{unc$\to$con} & \textbf{(\% gold)} & \textbf{(pp)} & \textbf{(\%)} \\
\midrule
Qwen3-8B  & SciERC & 1.95$\to$0    & 4 (0.41\%) & $-$0.10 & 0.22 \\
Llama-8B  & SciERC & 1.11$\to$0.09 & 5 (0.51\%) & $\pm$0.00 & 0 \\
Qwen3-8B  & DocRED & 1.03$\to$0.01 & 1 (0.01\%) & $-$0.84 & 2.05 \\
Llama-8B  & DocRED & 0.64$\to$0.21 & 1 (0.01\%) & $-$0.10 & 0.86 \\
\midrule
\multicolumn{2}{@{}l}{\textit{Average}} & & \textit{0.24\%} & $-$0.26 & 0.78 \\
\bottomrule
\end{tabular}}
\caption{Lucky violation analysis across four LoRA fine-tuned configurations.
Lucky: count (and \% of gold triples) of constraint-removed violations matching gold labels.
$\Delta$Recall: recall change from unconstrained to constrained.
Churn: fraction of gold-matching triples whose predictions change (including recovery via legal type paths).}
\label{tab:lucky-violations}
\end{table}

%%% Trie-Only Baseline %%%

\paragraph{Trie-Only Baseline (GCD-Equivalent).}
\label{app:trie-baseline}
Table~\ref{tab:trie-baseline} isolates the effect of vocabulary membership constraints without JSON structural enforcement.
The trie constrains entity type and relation tokens to schema-valid values (achieving ET${=}$0\%, RV${=}$0\% across all configurations), functionally equivalent to GCD's~\cite{geng2023gcd} constrained decoding.
However, without grammar-level JSON enforcement, models systematically omit entity type fields (83\% of triples) or produce malformed JSON (0--51\% parse failure), causing DR to rise to 15--96\%---far worse than unconstrained baselines.
The effect is most dramatic for Llama-8B LoRA: 50.8\% parse failure causes R-F1 to collapse from 45.47 (OntoCoDec S1) to 19.95, demonstrating that the grammar layer is not merely helpful but a structural prerequisite---vocabulary constraints require JSON scaffolding to reach the positions where they apply.

\begin{table}[h]
\centering
\small
\resizebox{\columnwidth}{!}{%
\begin{tabular}{@{}ll rrrr rr@{}}
\toprule
\textbf{Model} & \textbf{Setting}
  & \textbf{ET\,$\downarrow$} & \textbf{RV\,$\downarrow$} & \textbf{DR\,$\downarrow$} & \textbf{CF}
  & \textbf{R-F1\,$\uparrow$} & \textbf{PF\,(\%)} \\
\midrule
\multirow{2}{*}{Qwen3-8B}
  & LoRA  & 0 & 0 & 95.5  & 0.09 & 41.16 & 9.1 \\
  & ZS    & 0 & 0 & 29.3  & 0    & 8.64  & 0 \\
\multirow{2}{*}{Llama-8B}
  & LoRA  & 0 & 0 & 15.2$^\dagger$  & 0    & 19.95 & 50.8 \\
  & ZS    & 0 & 0 & 84.4  & 0.36 & 2.90  & 0 \\
\midrule
\multicolumn{8}{@{}l}{\textit{Comparison with OntoCoDec S1:}} \\
\midrule
& & \multicolumn{2}{c}{\textbf{S1 DR\,(\%)}} & \textbf{Trie DR\,(\%)} & \multicolumn{2}{c}{\textbf{DR Gap (pp)}} & \\
\cmidrule(lr){3-4} \cmidrule(lr){5-5} \cmidrule(lr){6-7}
Qwen3-8B  & LoRA  & \multicolumn{2}{c}{0}    & 95.5  & \multicolumn{2}{c}{95.5}  & \\
Qwen3-8B  & ZS    & \multicolumn{2}{c}{0.93} & 29.3  & \multicolumn{2}{c}{28.4}  & \\
Llama-8B  & LoRA  & \multicolumn{2}{c}{0.09} & 15.2  & \multicolumn{2}{c}{15.1}  & \\
Llama-8B  & ZS    & \multicolumn{2}{c}{0.08} & 84.4  & \multicolumn{2}{c}{84.3}  & \\
\bottomrule
\end{tabular}}
\caption{Trie-only baseline (GCD-equivalent, vocabulary membership without JSON grammar) on SciERC.
Top: full trie-only metrics (ET/RV/DR/CF: violation rates \%; R-F1 $\times$100; PF: parse failure rate).
Trie achieves ET${=}$0\%, RV${=}$0\% but DR rises to 15--96\% because models omit type fields without structural enforcement.
Bottom: DR gap between OntoCoDec S1 and trie-only, quantifying the grammar layer's contribution.
$^\dagger$50.8\% parse failure rate.}
\label{tab:trie-baseline}
\end{table}

%%% GCD/BoostCD Equivalence Comparison %%%

\paragraph{Comparison with Grammar-Based Constrained Decoding.}
\label{app:gcd-comparison}

The trie-only results above isolate vocabulary constraints in the absence of structural enforcement.
Table~\ref{tab:gcd-comparison} places this finding in the broader context of grammar-based constrained decoding methods, comparing the constraint class, enforcement level, and DR reduction achieved by each approach---including OntoCoDec's grammar layer, which is analytically equivalent to GCD-family methods.

\begin{table*}[t]
\centering
\small
\begin{tabular}{@{}l l l l l l@{}}
\toprule
\textbf{Method} & \textbf{Constraint Type} & \textbf{Level} & \textbf{DR Red.} & \textbf{CF} & \textbf{State} \\
\midrule
GCD~\cite{geng2023gcd} & CFG (JSON schema) & Grammar & ${\sim}$3\% & N/A & Stateless \\
BoostCD~\cite{boostcd2025} & CFG + beam rerank & Grammar & ${\sim}$3\%$^\dagger$ & N/A & Stateless \\
Trie-Only (Tab.~\ref{tab:trie-baseline}) & Prefix trie & Vocab & 0\%$^\ddagger$ & N/A & Stateless \\
\midrule
vLLM \texttt{guided\_json}~\cite{kwon2023vllm} & JSON Schema & Grammar & ${\sim}$3\%$^\S$ & N/A & Stateless \\
\midrule
OntoCoDec Grammar & CFG (XGrammar) & Grammar & ${\sim}$3\% & N/A & Stateless \\
OntoCoDec S1 & CFG + type constr. & Gram.+Sem. & 94--100\% (FT) & ${\leq}$1.4\% (FT) & Entity reg. \\
\bottomrule
\end{tabular}
\caption{Comparison with grammar-based constrained decoding methods.
$^\dagger$Equivalent mechanism; BoostCD adds beam-level reranking but the CFG enforcement is identical.
$^\ddagger$Vocabulary-only constraints without JSON structure cause DR to \emph{increase} (15--96\%; Table~\ref{tab:trie-baseline}).
$^\S$Experimentally verified: DR${=}$15.26\% vs.\ grammar-only 14.84\% ($\Delta{<}$0.5\,pp; Table~\ref{tab:vllm-verification}).
Grammar-only results are analytically equivalent to GCD/BoostCD (same CFG enforcement; \S\ref{sec:related_work}).}
\label{tab:gcd-comparison}
\end{table*}

Our grammar layer implements the same CFG-based token masking as GCD and BoostCD---the mechanisms are formally identical (JSON schema $\to$ pushdown automaton $\to$ per-step allowed token set).
Table~\ref{tab:scierc-ablation} row~2 therefore constitutes a direct comparison: grammar-only enforcement achieves RV/ET$\to$0\% but only ${\sim}$3\% DR reduction, identical to what GCD-family methods would achieve on IE tasks.

\paragraph{External verification via vLLM.}
To confirm this equivalence experimentally, we evaluated vLLM's \texttt{guided\_json} mode~\cite{kwon2023vllm}---an independent CFG enforcement implementation---on SciERC under identical conditions (Qwen3-8B, zero-shot, 551 documents).

\begin{table}[h]
\centering
\small
\resizebox{\columnwidth}{!}{%
\begin{tabular}{@{}l rrrr r@{}}
\toprule
\textbf{Method} & \textbf{DR\,$\downarrow$} & \textbf{RV\,$\downarrow$} & \textbf{ET\,$\downarrow$} & \textbf{CF} & \textbf{R-F1\,$\uparrow$} \\
\midrule
Grammar-Only (XGrammar) & 14.84 & 0 & 0 & 9.90 & 8.23 \\
vLLM \texttt{guided\_json} & 15.26 & 0 & 0 & --- & 8.75 \\
\midrule
$\Delta$ & $+$0.42 & 0 & 0 & --- & $+$0.52 \\
\bottomrule
\end{tabular}}
\caption{GCD-equivalence verification: OntoCoDec grammar-only vs.\ vLLM \texttt{guided\_json} on SciERC (Qwen3-8B, zero-shot, 551 documents). All differences $<$0.5\,pp, confirming that CFG-based approaches---regardless of implementation---leave ${\sim}$15\% DR violations unaddressed.}
\label{tab:vllm-verification}
\end{table}

All metrics differ by $<$0.5\,pp (Table~\ref{tab:vllm-verification}), confirming that any CFG-based constrained decoding tool---whether GCD, BoostCD, or vLLM---enforces the same constraint class and leaves the same ${\sim}$15\% DR gap.
The critical gap these methods cannot address is cross-field type consistency (Level~2), which requires the stateful semantic layer that maintains an entity registry across relation slots and enforces per-relation domain/range constraints at each type-bearing position.
